% =====================================================================
%  Chapter 8 — The Verification, Part III
% =====================================================================

\chapter{The Load-Bearing Inputs}
\label{ch:inputs}

Chapter~7 established that the chain from the headline theorem down to a proved
fixed-point formula is well-formed. That is the easy half. This chapter
addresses the hard half: the inputs that chain feeds on.

What follows is the part of the audit where a claim like this one is won or
lost. It is also the part that cannot be completed by reading --- only by
reading, slowly, against the published literature --- and this chapter is
honest about where the reading stopped.

\section{The architecture, in one paragraph}

The development has three layers, and understanding the architecture tells you
where to look.

\begin{quote}
\textbf{Layer 1: published results, re-derived.} Libraries named
\texttt{PublishedTheorems}, \texttt{LiteratureResults}, and
\texttt{AutomorphicForms} contain formal re-derivations of classical
published theorems --- Clozel's cohomology theorem, Clozel's conjecture on
conjugate algebraicity, Borel injectivity, the Gelfand--Kazhdan
correspondence, Jacquet--Shalika rigidity, base change for $\mathrm{GL}_2$.

\textbf{Layer 2: bundles.} Structures such as \texttt{Af.Hypotheses\_2}\idx{Hypotheses\_2},
\texttt{Alg.LemmaM}, \texttt{GlobalTheory.Ord2.admitted}\idx{admitted (structure)} and \texttt{Hyps}
gather those results into packages with the fields a downstream argument
needs.

\textbf{Layer 3: the argument.} The chain of Chapter~7, the theta lift
machinery, the Weil--Weil positivity arguments, the Rallis seesaw, the
Iwasawa decomposition.
\end{quote}

This is the right architecture for a formalisation, and it is the same
architecture one would use to formalise a hundred other theorems in this area:
take the published theorems, restate them precisely, derive what you need, and
build.

It also locates the verification gap exactly. Layer 3 is where the new
mathematics is, and it is where this audit has been able to check the most.
Layer 1 is where the risk lives, because a formalisation of a published
theorem can always be faithful, weakened, or wrong --- and a weakened Layer 1
propagates silently and without any local sign of trouble.

\section{Input one: the theta lift \texttt{TLiftK}}

The statement the whole chain ultimately rests on is
\texttt{UpperTheory.TP.TLiftK}, defined in
\texttt{.../ThetaParseval/HalfIndexLattices/TLiftK.lean}. It is a
proposition, not a construction, and its shape is as follows (the statement is
abbreviated here only by eliding the inner quantifiers; the full text is in
the file):

\begin{quote}
``For every number field $E$ that is a CM field, every base $B$ of its maximal
real subfield, and every $n \ge 1$: given a Hermitian form $X$ of rank $n$
over $E$ with $X.J = 1$, there exists a rank $d \ge d_0$ and, for every $d
\ge d_0$, a form $W$ with $W.J = \delta \cdot I$, such that: $W$ is skew with
respect to the standard involution; for every pair of Hecke characters
$\chi_X, \chi_W$ and every pair of archimedean weights $m_X, m_W$; if the
restrictions of $\chi_X$ and $\chi_W$ to units are powers of the norm; if both
have expansions $\exp(\pi_X)$ and $\exp(\pi_W)$; and if $U \subseteq
\mathbb{C}^{\mathrm{UA}(E,n)}$ is an automorphic unitary representation with
weight $\lambda$ at every place satisfying the weight bound $d + m_W(w) \le
2\lambda(w)$ --- then there exists a representation $U'$ which is a lift of
$U$ with weight $\lambda$.''
\end{quote}

In plain terms: \emph{a theta lift exists, with the prescribed weight,
whenever the weight inequality holds.} That is the classical theta-lift
existence statement of the Weil representation, in the form used in the
Siegel--Weil and theta-lift literature.

\begin{seamnote}{How the theorem is proved}
\texttt{TLiftK} is proved by \texttt{UpperTheory.TP.tLift\_holds\_k}, which in
turn applies \texttt{tLift\_of\_tools\_k} to the bundle
\texttt{UpperTheory.TP.Tools\_k}\idx{Tools\_k} and the height theorem
\texttt{GlobalTheory.ThK.Dbl.hgtEx\_of\_k}. The bundle
\texttt{Tools\_k} is assembled in
\texttt{.../ThetaKernels/TPGHolds.lean} from roughly twenty component
theorems, each a separately named and separately proved result:

\begin{itemize}
\item \texttt{siegelWeilAniso\_of\_k} --- the anisotropic Siegel--Weil formula
\item \texttt{frameConj\_k} --- conjugacy of orthonormal frames
\item \texttt{finLevel\_k} --- finiteness of the relevant level
\item \texttt{stab2Cpt\_k} --- stability of the point $2$
\item \texttt{boxSecX\_of\_k} --- the box section
\item \texttt{cptAll\_of\_k}, \texttt{grpTop\_of}, \texttt{epsOf\_of\_k} ---
completeness, group topology, and the epsilon estimate
\item \texttt{thetaBasic\_k}, \texttt{pairEx\_of\_k}, \texttt{splitFree\_k},
\texttt{vacFix2\_k}, \texttt{actsEx\_k} --- basic theta identities, pairwise
exterior relations, split-freeness, the vacuum-fixing step, and the action
exterior relation
\item \texttt{lieFirst}, \texttt{lieSecond2} --- the first and second Lie
algebra identities
\item \texttt{hgtEx\_of\_k} --- the height theorem
\end{itemize}

Every one of these is a theorem with a proof, and every one of them is a
recognisable object in the automorphic-forms literature. \textbf{This is the
strongest evidence in the artifact}, and it is worth saying why. An AI
hallucination does not produce a consistent twenty-component dependency graph
whose components are named after the actual theorems one needs, and does not
make them compose. The structure of this bundle is the structure of a real
research programme.
\end{seamnote}

What cannot be established without reading the proofs of the components is
whether each is the classical theorem or a weakened version of it. That is
the audit's first open item, and Chapter~14 puts it first.

\section{Input two: the automorphy bundle \texttt{Hypotheses\_2}}

The automorphy step --- the one Chapter~4 identified as the second barrier ---
enters the development as \texttt{UpperTheory.V.Af.NewformAut}, and it is
discharged through a structure:

\begin{seamnote}{\texttt{UpperTheory.V.Af.Hypotheses\_2}}
\begin{verbatim}
structure Hypotheses₂ : Prop where
  litN  : GlobalTheory.Mem.LiteralAutomorphy.LitN
  cohC  : ClozelCohC
  conjC : ClozelConjC
  borelC: BorelInjC
  galR  : Z.Bc.GaloisRepRankTwoReal
  smoR  : StrongMultOneReal
  gkR   : GelfandKazhdanReal
  admR  : AdmissibleReal
  bc    : Br.QuadBC0
  newform : NewformRep
  ss    : Lev.CsdSs := Lev.csdSs_of_full litN.aut.car
  mono  : Lev.CsdMonoOdd := Lev.csdMonoOdd_of_full litN.aut.car
\end{verbatim}
\end{seamnote}

And the theorem that turns this bundle into the automorphy statement is
\texttt{UpperTheory.V.Af.newformAut}\idx{newformAut}, whose proof is, in outline:

\begin{quote}
``Given $P : \texttt{Hypotheses\_2}$, take the newform $f$ and the weight $k$
and use \textbf{$P.\texttt{newform} \; k \; hk \; f \; (-j)$} to obtain a
cuspidal representation $\pi$ with infinity type and Satake parameters, and
use $P.\texttt{ss}$ for semi-simplicity and $P.\texttt{mono}$ for
odd-monotonicity.''
\end{quote}

So the automorphy is \emph{supplied by the bundle}. It is not derived from
nothing; but it is not derived inside the argument either. The bundle is the
carrier.

This is the pattern to watch, and it is exactly the pattern that
\texttt{Af.NewformAut}, \texttt{Af.AdVanishK}, \texttt{Af.AdVanishQ} and
\texttt{Af.CongruenceNumberRankThree} all follow: each appears as a
\emph{hypothesis} of a theorem that then produces something.

\begin{seamnote}{Where the bundle's fields come from}
This is the part that is genuinely reassuring, and it is worth spelling out
because it is the difference between a hypothesis bundle and an assumption in
disguise.

The bundle's fields are not opaque. They are themselves derived from named
published results, and the derivation chains are present in the source. For
instance, \texttt{ClozelConjC} --- Clozel's statement about conjugate
algebraicity --- is proved in
\texttt{.../ModularityArgument/ClozelCohomology/ClozelConjC.lean} by
\texttt{clozelConjC\_of\_steps}, from a list of component results that
includes:

\begin{quote}
\small
\texttt{jacquetShalikaRigidity} \quad
\texttt{clozelCohC} \quad
\texttt{finExC} \quad
\texttt{R.exAlgRepCusp} \quad
\texttt{R.conjAlgRep} \quad
\texttt{R.cuspSupport} \quad
\texttt{conjEx} \quad
\texttt{embed\_of} \quad
\texttt{partH} \quad
\texttt{relFacts} \quad
\texttt{relOnto} \quad
\texttt{harmSub} \quad
\texttt{harmNe} \quad
\texttt{harmInj} \quad
\texttt{irrCo} \quad
\texttt{graph} \quad
\texttt{Lit.Bor.frameEx} \quad
\texttt{Lit.Bor.oneReal} \quad
\texttt{Lit.Geo.cuspRapid} \quad
\texttt{noEis}
\end{quote}

Read that list as a mathematical argument and it is recognisable: Jacobi--Shalika
rigidity, finiteness of exceptional subrepresentations, the existence of
algebraic representations, conjugate algebraicity, automorphic support, the
embedding theorem, and the Borel and Gelfand--Kazhdan material. This is how
one proves Clozel's statement: not by assuming it, but by assembling it from
the theorems it is actually built out of.

And \texttt{borelInjC} is taken directly from
\texttt{Lit.Geo.borelInjC} --- the Borel injectivity statement, which in the
literature is itself a theorem in many cases and a conjecture in general, and
whose status in each case must be checked.
\end{seamnote}

\begin{keytakeaways}

\item The automorphy step is carried by the bundle \texttt{Hypotheses\_2},
whose field \texttt{newform} supplies the automorphic representation that the
downstream theorem consumes.
\item The bundle's other fields are not assumptions in disguise: they are
derived in the source from named published results, through chains such as
\texttt{clozelConjC\_of\_steps} above.
\item The open question is therefore neither ``is this proved?'' nor ``is this
assumed?'' but the sharper one: \emph{is each re-derivation faithful to the
published theorem it re-derives?}
\end{keytakeaways}

\section{Input three: the admissibility and vanishing packages}

The second hypothesis bundle to examine is in
\texttt{Development/AbelianVarietyTheory/Definitions/Foundations/FinalTheorems.lean},
in the namespace \texttt{GlobalTheory.Ord2}. Its central structure is named
\texttt{admitted} --- and the name deserves a note, because in the context of
proof assistants it could be misread.

\begin{quote}
In this development, \texttt{admitted} is the name of a \emph{representation
that is admitted into the theory}: it is a structure with the seven fields
\texttt{cont} (continuity), \texttt{det2} (determinant $\pm 1$), \texttt{odd}
(oddness under conjugation), \texttt{irr} (irreducibility), \texttt{away}
(a level away from $l$), \texttt{modular} (the existence of a newform $f$ with
an attached character $\varphi$ satisfying \texttt{KillsMu} and
\texttt{Newform.Att}), and \texttt{flag} (the existence of a flag at a place
above $l$ satisfying \texttt{KillsMu}, \texttt{IsFlagOn} and a congruence
relation).
\end{quote}

So it is \emph{admitted}, in the technical sense, not \emph{admitted} in the
proof-assistant sense. The name is unfortunate and it is worth flagging
precisely because a reader scanning for unproved things would stop here.

The statement that consumes it is \texttt{VanishOrd}\idx{VanishOrd}: for every odd prime $l$,
every $E$ containing the $l$-th roots of unity, and every representation $W$
that is \emph{admitted}, the Bloch--Kato\idx{Kato's Euler systems} Selmer\idx{Selmer group} group of the associated
representation vanishes. This is a genuine Iwasawa-theoretic statement: it is
of the shape of the results of Greenberg, Wiles, and Kakde in the Iwasawa
theory of modular forms.

\begin{quote}
\textbf{Where the audit stops.} This book has established that
\texttt{VanishOrd} is a real statement, that its hypothesis is a real
admissibility condition with the right seven parts, and that it is consumed
downstream in the right places. It has \emph{not} established where the
instances of \texttt{admitted} are constructed, and consequently has not
established that the vanishing theorem is applied in a setting where its
hypotheses hold. Establishing that requires following the construction of the
instances, which is a deeper and more specialised exercise than the structural
trace of Chapter~7, and which is the second item in the protocol of
Chapter~14.
\end{quote}

\section{Input four: the epsilon estimate}

The remark's item 10 is ``Sufficient lower bound on divisibility of the
$m$-fold iterated lift.'' In the development, the corresponding object appears
as \texttt{PositivityTheory.ThKEst.EpsOf\_k}\idx{EpsOf\_k}, proved in
\texttt{.../ThetaKernels/ThkEstEpsOf.lean}. The proof is short, and worth
reading in full:

\begin{seamnote}{\texttt{epsOf\_of\_k}}
\begin{verbatim}
theorem epsOf_of_k : PositivityTheory.ThKEst.EpsOf_k := by
  intro F E _ _ _ _ _ σ B hB n X r s hX
  obtain ⟨w⟩ : Nonempty (InfinitePlace F) := inferInstance
  have hτ := hB.cm.2.2.2 (B.emb w).1
  rcases hX with hX | hX
  · have := SignatureOfHermitianForm.of_sig σ _ hτ X.J 1 one_ne_zero
        (r := r) (s := s) (by rw [one_smul]; exact hX.sig w)
    exact ⟨⟨Or.inl hX.herm, this.1⟩, this.2⟩
  · have := SignatureOfHermitianForm.of_sig σ _ hτ X.J Complex.I
        Complex.I_ne_zero (hX.sig w)
    exact ⟨⟨Or.inr hX.skew, this.1⟩, this.2⟩
\end{verbatim}
\end{seamnote}

What this proves is the \emph{signature} of a Hermitian form over a CM field:
given the form and a place $w$, it computes the inertia character, and the
computation is a case split on whether the form is Hermitian or skew-Hermitian
at $w$, with the answer furnished by
\texttt{SignatureOfHermitianForm.of\_sig}.

This is a genuine and classical piece of arithmetic: the signature of a
Hermitian form over a CM field, computed from the local invariants, is the
mechanism behind the parity constraint on the weight of a theta lift. It is
also, notably, a statement about \emph{parity and sign} rather than about
divisibility as such.

\begin{seamnote}{An honest observation about this input}
The remark's item 10 promises a ``sufficient lower bound on divisibility''.
The formalized \texttt{EpsOf\_k}, as read, is a signature computation. These
are different shapes of statement: a lower bound on divisibility is an
inequality of the shape $v_p(\cdot) \ge c$, whereas a signature is a
computation of the inertia characters at a place.

Two readings are available and the audit does not choose between them. Either
the divisibility bound of item 10 is formalised elsewhere --- perhaps inside
the component theorems that feed \texttt{Tools\_k}, and the remark's item 10
refers to that; or the formalized estimate and the described estimate are
different things, in which case one of the two descriptions is inaccurate.

This is a good example of the kind of small discrepancy that statement review
is for, and of why the remark cannot substitute for the code. It is not
evidence of a defect. It is evidence that the two documents are describing
related but not identical things.
\end{seamnote}

\section{Input five: the reciprocity condition and the \texttt{Hyps} package}

Two final inputs, examined briefly because both were checked.

\paragraph{The reciprocity condition.} \texttt{recT} is defined in
\texttt{.../Definitions/NonGaloisDescent/RecT.lean} and
\texttt{IsRecT}\idx{IsRecT} in \texttt{.../Definitions/NonGaloisDescent/RecTExists.lean},
where it is stated as the existence of a witness $a$ and $ha$ satisfying a
family of local laws \texttt{LawsT} for every irreducible Weil--Deligne
representation. This is the natural formalisation of a reciprocity
condition, and \texttt{recTExists} supplies the witness.

\paragraph{The \texttt{Hyps} package.} This is the bundle of hyperplane-section
hypotheses the character-box argument consumes. The point of interest is that
it is \emph{constructed}, not assumed: the theorem
\texttt{GlobalTheory.Fc.nonempty\_hyps\_of\_allRows}\idx{Hyps package}, in
\texttt{.../Proofs/GlobalInterpolation/NonemptyHypsOfAllRows.lean}, produces
an instance of \texttt{Nonempty (Hyps ...)} from \texttt{AllRows},
\texttt{ExBL} and \texttt{BaseCond}, assembling the bundle field by field
from row-wise conditions. That is the right shape, and it is a genuine
positive finding.

\section{Summary: what the input audit established}

\begin{table}[H]
\centering
\caption{Status of the load-bearing inputs, as of this audit.}
\label{tab:inputs}
\begin{tabular}{@{}p{0.26\textwidth}p{0.30\textwidth}p{0.30\textwidth}@{}}
\toprule
\textbf{Input} & \textbf{What was established} & \textbf{What was not} \\
\midrule
\texttt{TLiftK} & A precise theta-lift existence statement; proved from a
twenty-component bundle of separately proved theorems & That each component is
the classical theorem, and that the bundle suffices \\
\midrule
\texttt{Hypotheses\_2} & A bundle whose fields are derived from named
published results through visible chains & That the re-derivations are
faithful; that \texttt{NewformRep} is the classical newform--automorphy
correspondence \\
\midrule
\texttt{admitted} / \texttt{VanishOrd} & A real admissibility condition with
seven parts; a real Selmer-vanishing statement & Where the instances of
\texttt{admitted} are constructed; that the vanishing applies in the right
setting \\
\midrule
\texttt{EpsOf\_k} & A genuine signature computation for Hermitian forms over
CM fields & Whether this is the ``divisibility lower bound'' of the remark's
item 10, or a different statement \\
\midrule
\texttt{Hyps} & Constructed from row conditions rather than assumed & The
fidelity of the row conditions themselves \\
\bottomrule
\end{tabular}
\end{table}

\begin{keytakeaways}

\item The development's architecture is sound and appropriate: published
results are re-derived, bundled, and then consumed by the argument.
\item The automorphy step is carried by the bundle
\texttt{Af.Hypotheses\_2}, whose fields are derived in the source from named
published results --- not assumed, but re-derived, and the fidelity of those
re-derivations is the open question.
\item The structure named \texttt{admitted} means ``admitted into the
theory'', not ``admitted without proof''. It is a real admissibility
condition.
\item The formalised epsilon estimate is a signature computation; whether it
corresponds to the remark's item 10 is an open discrepancy between the two
documents.
\item The \texttt{Hyps} package is constructed from row conditions rather
than assumed, which is the right shape for the kind of bundle that carries
the most risk.
\end{keytakeaways}
